\documentclass[10pt]{article}

% Specify the margins and text size.
\setlength{\textwidth}{6.4in}
\setlength{\textheight}{9.5in}
\setlength{\oddsidemargin}{0pt}
\setlength{\evensidemargin}{0pt}
\setlength{\topmargin}{0pt}
\setlength{\hoffset}{.05in}
\setlength{\voffset}{-1in}

\setlength{\parskip}{5pt}
\setlength{\parindent}{0pt}

% Load some fonts and symbol packages
\usepackage{latexsym}
\usepackage{pifont}       % contains 'star' symbol for counterinsurgency handbook title
\usepackage{yfonts} 
\usepackage{amsmath}
\usepackage{amsfonts}

\usepackage{graphicx}     % actually, this is loaded by pstricks
\usepackage[T1]{fontenc}
\usepackage{ifthen}
\usepackage{pstricks,pst-grad,pst-text,pst-node,multido,pst-plot,pst-func,calc,pst-3dplot}
\usepackage[all]{xy}
%\usepackage{animate}

% The hyperref package inserts links into the pdf.
\definecolor{MyLinkColor}{rgb}{.1,.2,1}
\definecolor{MyCiteColor}{rgb}{.1,1,.2}
\definecolor{MyURLColor}{rgb}{.4,.4,.4}
\usepackage[backref=true,pagebackref=false,hyperindex,colorlinks=true,
  linkcolor=MyLinkColor,urlcolor=MyURLColor]{hyperref}


% The tweaklist package is something I found on the web.  It provides a simple interface
% for making changes to spacing used in the itemize and enumerate environments.  Comment
% this out if you don't care to use tweaklist.
\usepackage{tweaklist}
\renewcommand{\itemhook}{\setlength{\parskip}{2pt}\setlength{\parsep}%
{1pt}\setlength{\topsep}{0pt}\setlength{\itemsep}{0pt}}

\newcommand{\U}{\underline{\hspace{5pt}}}

\usepackage{listings}
\newcommand{\Z}{\hphantom{0}}

\begin{document}
\pagestyle{empty}
\lstset{language=R, showspaces=false, showstringspaces=false}

\href{http://www.su.edu}{\includegraphics[height=1.75cm]{sulogo.eps}}
\vspace{-1.69cm}
\bigskip

\setlength{\baselineskip}{1.05\baselineskip}

\begin{center}
\textbf{\large  Exam IV}
\end{center}

1. (25 points) A survey is conducted of 100 freshmen at public universities.
Among the surveyed  students, 20\% (or 1/5) work at least 15 hours per week for pay.

\hspace{20pt} a) Find a 95\% confidence interval
on the  percent of freshmen at public universities who work at least 15 hours per week for pay.
\vspace{2.3in}

\hspace{20pt} b) True or false and explain:  There is a 95\% chance that the
percent among \textit{all} freshmen at public universities is in the range you found in a).
\vspace{1.5in}

\hspace{20pt} c) True or false and explain:  There is a 95\% chance that
if a second sample of 100 freshmen were taken,
the sample percent would be in the range you found in a).
\vspace{1.5in}

\hspace{20pt} d) True or false and explain:  There is a 95\% chance that the
average hours a week worked for pay among all 18--20 year olds
is in the range you found in a).
\vspace{1in}
\vfill
\eject
{\ }



2. (10 points)
Would taking the average of 100 measurements divide
the likely size of the chance error by a factor of 5, 10 or 25? Justify your answer.
\vspace{1.8in}

3. (10 points)

\hspace{10pt} a) Other things being equal, which of the following $P$-values is best for the null
hypothesis?  Explain.
\[0.1\% \hspace{.5in} 3\%\hspace{.5in} 17\%\hspace{.5in} 32\%\]
\vspace{.4in}

\hspace{10pt} b) Repeat b) for the alternative hypothesis.
\vspace{.9in}

4.  (15 points) The Zorro News Organization reports that  50\%
of freshmen at public universities  work 15 hours a week for pay.
If that is report is accurate, approximately what is the chance of getting a sample
percent as low as 20\% (as we did in \#1) or lower?  Hint:  Use a $z$-test.
\vfill
\eject


5. (15 points) Five hundred draws are made at random from the box
\begin{center}
\begin{pspicture}(0,0.1)(4,0.3)
\psline(0,0.5)(0,0)(4,0)(4,0.5)
\rput[l](0.15,0.3){60,000\ \framebox{0}\,s}
\rput[l](2.15,0.3){20,000\ \framebox{1}\,s}
\end{pspicture}
\end{center}
Note that $\sqrt{0.25\times 0.75}\,/\sqrt{500} \approx 0.02$. True or false?  Explain your conclusions.
\smallskip

\begin{itemize}
\item[i.] The expected value for the percentage of 1s among the draws is exactly 25\%.\\[.5in]
\item[ii.] The expected value for the percentage of 1s among the draws is around 25\%, give
or take 2\% or so.\\[.5in]
\item[iii.] The percentage of 1s among the draws will be around 25\%, give or take 2\% or so.\\[.5in]
\item[iv.] The percentage of 1s among the draws will be exactly 25\%.\\[.5in]
\item[v.] The percentage of 1s in the box is exactly 25\%.\\[.5in]
\item[vi.] The percentage of 1s in the box is around 25\%, give or take 2\% or so.
\end{itemize}
\vspace{1in}

6. (10 points) Find the specified area under the normal curve.  Write down the \texttt{R} command that you use.

\hspace{10pt} a) $z<-0.35$ 
\bigskip
\bigskip\bigskip\bigskip

\hspace{10pt} b) $z>-1.35$
\bigskip
\bigskip\bigskip\bigskip

\hspace{10pt} c) $-1.5 < z < 0.8$
\bigskip\bigskip\bigskip\bigskip


\vfill
\eject

\newcommand{\BL}{\underbar{\hspace{.5in}}}
7. (15 points)
Laser altimeters can measure elevation to within a few inches, witut bias,
and with no trend or pattern to the measurements.
As part of an experiment, 25 readings were made on the elevation of a mountain peak.
These averaged out to 81,411 inches, and their SD was 30 inches. 
\bigskip
\bigskip

\hspace{20pt} a) The elevation of the mountain peak is estimated 
as \BL\BL; this estimate is likely to be\vspace{10pt}

\hspace{20pt}\hphantom{f) } off by \BL\BL\ or so.
\vspace{1.0in}

\hspace{20pt} b) (T/F) 81,411$\pm$ 12 inches is a 95\%-confidence interval for the elevation
of the mountain peak.
\vspace{1.0in}

\hspace{20pt} c) (T/F) 81,411$\pm$ 12 inches is a 95\%-confidence interval for 
average of the 25 readings.
\vspace{1.0in}

\hspace{20pt} d) (T/F) There is about a 95\% chance that the next reading will be in 
the range 81,411 $\pm$ 12 inches.
\vspace{1.0in}

\hspace{20pt} e) (T/F) About 95\% of the readings were in the range 81,411 $\pm$ 12 inches.
\vspace{1.0in}

\hspace{20pt} f) If another 25 readings are made, there is about a 95\% chance that their
average will be in the range\vspace{-4pt}

\hspace{20pt}\hphantom{f) } 81,411 $\pm$ 12 inches.


\vfill
\eject

\end{document}

